\section{Introduction}\label{sec:introduction}

Restricted geometric Langlands theory compares automorphic sheaves with
sheaves on a stack of local systems whose variation is restricted in the
sense of Arinkin--Gaitsgory--Kazhdan--Raskin--Rozenblyum--Varshavsky
\cite{AGKRRV}. In positive characteristic, Gaitsgory--Raskin construct an
equivalence onto an open-and-closed union of spectral components
\cite[Main Theorem~1.3.9(i)]{GR}. The remaining support question asks whether
any component can be absent. We prove that none is absent in the two
precise settings below.

The geometric comparison grew from the construction of individual Hecke
eigensheaves. Over a finite field, unramified automorphic functions can be
viewed as functions on isomorphism classes of bundles; Frobenius traces
turn sheaves on the bundle stack into such functions. For general linear
groups, the constructions of Drinfeld and Laumon led to the question of
attaching a Hecke eigensheaf to each irreducible local system.
Frenkel--Gaitsgory--Vilonen reduced this existence problem to a vanishing
conjecture, which they established over finite fields using Lafforgue's
work \cite[\S0.1 and \S\S1--2]{FrenkelGaitsgoryVilonen}.
Gaitsgory subsequently proved the vanishing conjecture geometrically
\cite[\S0.1 and \S2]{GaitsgoryVanishing}.
The categorical formulation of Beilinson--Drinfeld seeks a comparison
of the whole automorphic category with a category varying over the stack
of local systems \cite[\S0.1]{AGKRRV}.

The choice of spectral category is essential. In characteristic zero,
Arinkin--Gaitsgory explained why quasi-coherent sheaves on the local-system
stack do not suffice for a general reductive group, and formulated the
comparison using ind-coherent sheaves with nilpotent singular support
\cite[\S\S1.1.1--1.1.6]{ArinkinGaitsgory}.
The extra singular direction is a nilpotent horizontal section of the
adjoint local system, reflecting the additional datum in an Arthur parameter.
Their formulation also satisfies the compatibility tests supplied by
geometric Satake and Eisenstein series. On the automorphic side,
Beilinson's theory supplies singular support for constructible sheaves
in arbitrary characteristic \cite[\S1.3]{BeilinsonSingularSupport}.
These two support theories concern different categories and play
different roles in the correspondence.

To formulate a categorical comparison for \(\ell\)-adic sheaves,
Arinkin--Gaitsgory--Kazhdan--Raskin--Rozenblyum--Varshavsky introduced
the restricted local-system stack \cite[\S\S0.1--0.2 and \S0.5.4]{AGKRRV}.
Its families are defined through the tensor category of local systems
in the chosen sheaf theory, rather than through a de Rham moduli problem.
The connected components are indexed by semisimple local systems,
while each component retains extensions and derived deformation data
\cite[Proposition~3.7.2 and Corollary~3.7.4]{AGKRRV}.
They constructed the coherent spectral action on automorphic sheaves
with nilpotent singular support and proved, under their characteristic
hypotheses, that Hecke eigensheaves have nilpotent singular support,
as predicted by Laumon
\cite[\S0.2.3 and Theorems~14.3.2 and~14.4.3]{AGKRRV}.
This supplies both the categories and the Hecke-compatible action used
in the present paper.

The characteristic-zero geometric Langlands theorem was established in
the five-part project of Arinkin, Beraldo, Campbell, Chen, Faergeman,
Gaitsgory, Lin, Raskin, and Rozenblyum. Gaitsgory--Raskin's construction
of the functor and their multiplicity-one theorem are the first and
final parts of that project \cite{GLCI,GLCV}.
Their subsequent work \cite{GR} uses geometric specialization to relate
the characteristic-zero and positive-characteristic \(\ell\)-adic theories.
It proves the primed equivalence used here and the full result for
general linear groups. Our input from characteristic zero is precisely
\cite[Main Theorem~1.3.9(ii)]{GR}.

An equivalence onto a union of components does not yet give an
automorphic category at a parameter outside that union. The issue is
therefore existence on the omitted components, even after the
categorical comparison has been constructed. Our proof addresses
this issue by showing that the geometric specialization functor
cannot annihilate the nonzero spectral summand attached to such a
component. The local tests that establish this fact are described
after the main statements.

\subsection{The spectral stack and the two hypothesis regimes}

Let $k$ be an algebraically closed field of characteristic $p>0$, let
$\ell\ne p$ be a prime, and set $E=\overline{\mathbb Q}_\ell$.
Let $X/k$ be a smooth projective connected curve and $G/k$ a connected
reductive group. Write $\check G/E$ for the Langlands dual group and
$\mathfrak g=\Lie(G)$. Such a group $G$ is split over $k$.

We use the $E$-linear sheaf categories of \cite{AGKRRV,GR}.
In particular, $\Shv(X)$ is the ind-constructible category, and
$\QLisse_E(X)$ denotes the full dg subcategory of $\Shv(X)$ whose ordinary
cohomology sheaves are filtered colimits of finite-rank lisse $E$-sheaves
\cite[\S1.2.5 and Definition~1.2.6]{AGKRRV}. Its distinction from the
ind-completion $\IndLisse_E(X)$ of constructible lisse complexes will
matter in Section~\ref{sec:specialization}. The restricted spectral
prestack $Y=LS^{\mathrm{restr}}_{\check G}(X)$ is defined on derived affine
$E$-schemes $T$ by
\begin{equation}\label{intro:restricted-stack}
 Y(T)=\Fun^{\otimes,\mathrm{r.t.ex}}
 \bigl(\Rep_E(\check G),\QCoh(T)\otimes_E\QLisse_E(X)\bigr).
\end{equation}
Here the superscript denotes right-t-exact symmetric monoidal functors,
with the conventions of \cite{AGKRRV}. A parameter is an $E$-point of
$Y$, hence such a tensor functor with $T=\Spec E$.

Write $\Bun_G(X)$ for the stack of principal $G$-bundles on $X$.
Its global nilpotent cone consists of cotangent Higgs fields with
nilpotent values. It defines the full dg subcategory
$\Shv_{\Nilp}(\Bun_G(X))$. On $Y$ we use the spectral nilpotent
singular-support condition and the category $\IndCoh_{\Nilp}(Y)$ of
\cite[\S\S21.2.1--21.2.5]{AGKRRV}. These are the established categories; the
new assertion concerns the spectral components on which the equivalence
is supported.

We state the characteristic assumptions individually. A Levi subgroup
always means a Levi factor of a parabolic subgroup. Semisimple Lie
elements are geometrically conjugate into the Lie algebra of a maximal
torus; nilpotent elements have zero in their geometric adjoint-orbit
closure.

\begin{hypothesis}[Common characteristic assumptions]\label{hyp:common}
The following conditions hold for $G/k$.
\begin{enumerate}[label=\textup{(\roman*)}]
\item There is a nondegenerate $G$-invariant symmetric bilinear form
$\kappa$ on $\mathfrak g$ whose restriction to the center of
$\Lie(M)$ is nondegenerate for every Levi subgroup $M$.
\item For every Levi subgroup $M$, including $G$, and a maximal torus
$T_M\subset M$, Chevalley restriction is an isomorphism
\[
 k[\Lie(M)]^M\xrightarrow{\ \sim\ }
 k[\Lie(T_M)]^{W_M},\qquad W_M=N_M(T_M)/T_M.
\]
\item The scheme-theoretic centralizer in $G$ of every semisimple element of
$\mathfrak g$ is a Levi subgroup.
\item For every field extension $k'/k$, every nilpotent element of
$\mathfrak g\otimes_k k'$ belongs to $\Lie(R_u(P))$ for some
parabolic subgroup $P\subset G_{k'}$ defined over $k'$.
\end{enumerate}
\end{hypothesis}

The first three conditions are those of \cite[\S14.4.1]{AGKRRV}; the
last is the additional condition of \cite[\S D.1.1]{AGKRRV}.
These are the assumptions referenced by \cite[\S0.1.9]{GR}.
The center in condition~(i) is the center of a Lie algebra. We will
justify its relation to central cocharacters where that relation is used.

\begin{hypothesis}[Regime A: finite-field descent]\label{hyp:A}
There are $q=p^r$, a smooth projective geometrically connected curve
$X_0/\mathbb F_q$, and a split connected reductive group
$G_0/\mathbb F_q$ such that
\[
 k=\overline{\mathbb F}_q,\qquad
 X=X_0\times_{\mathbb F_q}k,\qquad
 G=G_0\times_{\mathbb F_q}k.
\]
Hypothesis~\ref{hyp:common} holds. No additional restriction on the
order of the Weyl group is imposed in this regime.
\end{hypothesis}

\begin{hypothesis}[Regime B: arbitrary algebraically closed base]
\label{hyp:B}
The field $k$ is any algebraically closed field of characteristic $p>0$.
Hypothesis~\ref{hyp:common} holds, $p$ is very good for $G$, and
$p\nmid |W_G|$. No finite-field descent of $X$ is assumed.
\end{hypothesis}

Here very good has its usual root-system meaning: for type $A_n$,
$p\nmid n+1$; for $B_n,C_n,D_n$, $p\ne2$; for
$G_2,F_4,E_6,E_7$, $p\ne2,3$; and for $E_8$, $p\ne2,3,5$,
on every irreducible factor. A torus has no such exclusions.
We retain all the conditions in Hypothesis~\ref{hyp:B}, even where
some are redundant for a particular step. The two regimes have distinct
proofs of the needed Lie-theoretic statement; neither is used as a
substitute for the other.

\subsection{Full support}

Under Hypothesis~\ref{hyp:common}, the restricted Langlands functor
of \cite{GR} factors as
\begin{equation}\label{intro:primed}
 \mathbb L_G^{\mathrm{restr}}:
 \Shv_{\Nilp}(\Bun_G(X))
 \xrightarrow{\ \sim\ }\IndCoh_{\Nilp}(Y')
 \hookrightarrow\IndCoh_{\Nilp}(Y),
\end{equation}
where $Y'\subset Y$ is an open-and-closed union of connected
components. The functor is $\QCoh(Y)$-linear
\cite[Lemma~1.3.7 and Main Theorem~1.3.9]{GR}.

\begin{theorem}[Full support]\label{thm:full-support}
Let $k,\ell,E,X,G$ be as above. Under either
Hypothesis~\ref{hyp:A} or Hypothesis~\ref{hyp:B}, the inclusion
$Y'\subset Y$ is equality as spectral prestacks. Consequently the
given functor is a $\QCoh(Y)$-linear equivalence
\[
 \mathbb L_G^{\mathrm{restr}}:
 \Shv_{\Nilp}(\Bun_G(X))\xrightarrow{\ \sim\ }
 \IndCoh_{\Nilp}(Y).
\]
\end{theorem}

This proves Conjecture~1.3.10 of \cite{GR} in the two displayed
regimes. The positive-characteristic primed equivalence, the full
characteristic-zero equivalence, and the full result for general linear
groups are established inputs from Main Theorem~1.3.9 of that paper.
The present theorem concerns the restricted theory with its stated
characteristic hypotheses.

Section~\ref{sec:consequences} records three consequences. Localizing
the given functor at any spectral component gives the corresponding
nonzero ind-coherent category. In regime A, the known primed trace
formula becomes the full arithmetic formula using derived geometric
Frobenius fixed points. For an arbitrary parameter, spectral linearity
produces a nonzero Hecke eigenobject with coherent tensor and fusion
compatibilities. Its asserted category is ind-constructible and
nilpotent; no bounded constructibility or perversity assertion is needed.

\subsection{Rational coefficients}

There is also a coefficient form of the support question. Put
$E_0=\mathbb Q_\ell$, retaining $E=\overline{\mathbb Q}_\ell$.
For the rational support statement let $k=\overline{\mathbb F}_q$, and let $X/k$ be
smooth projective connected and $G/k$ connected reductive satisfying
Hypothesis~\ref{hyp:common}. Let $\check G_0/E_0$ be the split dual group,
with $\check G=\check G_0\times_{E_0}E$. Define $Y_0$ by
\eqref{intro:restricted-stack} with $E_0$ in place of $E$, and put
$\mathcal C_0=\Shv_{\Nilp,E_0}(\Bun_G)$, using the rational
ind-constructible sheaf category and the same nilpotent condition.
These coefficient changes leave the geometric field $k$ unchanged.

Universal evaluation sends a representation to its local system in the
universal family on $Y_0$; for a finite set of legs it uses the exterior
products of these local systems. The rational spectral-action datum
consists of a continuous
$\QCoh(Y_0)$-action on $\mathcal C_0$ whose restriction along universal
evaluation is the usual multi-leg Hecke system, with its tensor,
permutation and collision compatibilities. We take this datum as part
of the rational support problem. The rational geometric Satake and
derived Hecke constructions are recalled in Section~\ref{sec:coefficients}.
We distinguish the given rational action and support from a construction
of a normalized rational Langlands equivalence: the cited theorem of
Gaitsgory--Raskin is used over $E$.

\begin{theorem}[Rational spectral support]\label{thm:rational-support}
For the rational datum just specified, every nonempty open-and-closed
substack $U\subset Y_0$ has
\[
 \QCoh(U)\underset{\QCoh(Y_0)}\otimes\mathcal C_0\ne0.
\]
Consequently, if $Y'_0\subset Y_0$ is an open-and-closed support union
whose complementary structure idempotent acts by zero on
$\mathcal C_0$, then $Y'_0=Y_0$
as spectral prestacks. In particular this applies to the support of
any given $\QCoh(Y_0)$-linear primed equivalence.
\end{theorem}

The conclusion concerns the entire rational spectral prestack, without
identifying its components with rational points. A curve over
$\overline{\mathbb F}_q$ descends to a finite extension of
$\mathbb F_q$ by finite presentation; a split form of $G$ does as well.
Thus Theorem~\ref{thm:full-support} in regime A supplies the geometric
input without adding a Weyl-order restriction. No rational extension
of regime B is asserted here.

Section~\ref{sec:coefficients} proves the required coefficient comparison
on sheaf categories, affine families and the coherent spectral action.
It then descends the absence of a clopen summand acting by zero.
The rational arithmetic statement has a further input: a canonical
primed trace formula over $E_0$ with the stated Frobenius and
ind-coherent conventions. We prove unpriming conditional on that input;
we do not construct the rational primed comparison map. The
unconditional arithmetic formula above remains over $E$.

\subsection{The specialization and slicing argument}

We first outline the proof of Theorem~\ref{thm:full-support} over $E$.
A missing spectral component leads to a contradiction as follows.
Lift $X$ to a smooth proper curve $\mathscr X$ over the Witt trait,
lift $G$, and let $\mathcal B$ be the relative bundle stack. Write
$K$ for an algebraic closure of the fraction field and $\Psi$ for
geometric nearby cycles with monodromy forgotten. The established
specialization formalism and characteristic-zero equivalence produce,
from a missing component, a nonzero compact sheaf $F$ on
$\mathcal B_K$ with $\Psi F=0$.

Section~\ref{sec:specialization} strengthens this vanishing to Hecke
transforms pulled back along maps whose bundle projection is smooth.
The leg map may be arbitrary. This extra freedom is needed after
cutting a Hecke correspondence by a hypersurface.

On a smooth chart $U\to\mathcal B$, a special-fiber cotangent vector
outside the nilpotent cone gives a test function $f$. If that covector
comes from a Higgs field, Sections~\ref{sec:lie} and~\ref{sec:hecke}
construct a modification with a nonzero leg covector and transverse
zeros on both fixed-side Hecke fibers. The two characteristic regimes
enter only in constructing the Levi and central cocharacter required
for this calculation.

Section~\ref{sec:twohecke} takes two modifications with common output
and common leg. Proper pushforward and the reverse transversality
isolate the stalk at the diagonal. The forward transversality then
gives an actual \'{e}tale local equation
\[
 \sum_{i=1}^m u_i v_i+f=0,
\]
with the sheaf pulled back from $U$. Here $m$ is the Hecke-orbit
dimension and $(u_i,v_i)$ are local coordinates transverse to the
chart and leg. This is the point where the two correspondences are
both essential.

Section~\ref{sec:slices} compares the projective closure of this
quadratic equation with its hyperplane classes. The remaining derived
sheaf is a constant line supported on $f=0$. A projective incidence
construction then gives the same vanishing on higher-codimension
slices whose cotangent planes avoid the nilpotent cone.
Finally, Section~\ref{sec:support} uses the nilpotent dimension bound
to choose such a slice through the actual support of $F$. Its support
is finite over a henselian trait near the chosen point, so its nearby
cycles are a finite direct sum of stalks with at least one nonzero
summand. This contradicts slice vanishing.

The central-in-a-Levi modification and moving-leg calculation originate
in \cite[\S\S20.4--20.7]{AGKRRV}. The additional argument here is the
two-correspondence transversality and slicing mechanism, including
the exact local equation and derived comparison triangles. All
singular-support dimension bounds used in the proof are on the
special fiber; the nearby-cycle tests themselves are proved directly.
The arithmetic application combines the automorphic trace comparison
of \cite{AGKRRVTrace} with the spectral trace calculation of
Beraldo--Lin--Reeves \cite[\S\S6.4.12--6.4.13]{BeraldoLinReeves}, as
assembled in \cite[\S1.5]{GR}. The coefficient argument in
Section~\ref{sec:coefficients} uses constructible coefficient
continuity \cite{HemoRicharzScholbach} and the universal-local-acyclicity
criterion of \cite{HansenScholze}; the rational spectral action and
primed arithmetic map remain the given data stated above.
